\newif\iffable
\fabletrue    % Uncomment to see instructions from claude fable

\chapter{RNS Gadgets}
% \ifnotes
%     Recall that:
%     \begin{equation*}
%     a = (c_0, c_1, \dots c_{N-1}) \pmod{Q} =
%     \begin{pmatrix}
%         c_0 \bmod{q_0} & c_1 \bmod{q_0} & \dots & c_{N-1}\bmod{q_0} \\
%         c_0 \bmod{q_1} & c_1 \bmod{q_1} & \dots & c_{N-1}\bmod{q_1} \\
%         \vdots & \vdots & \ddots & \vdots \\
%         c_0 \bmod{q_k} & c_1 \bmod{q_k} & \dots & c_{N-1}\bmod{q_k}
%     \end{pmatrix}
%     \end{equation*}
% \fi

\section{RNS Gadgets}
\iffable{
    \color{purple}
    \begin{itemize}
        \item Establish HPS, BV and GHS/Hybrid gadgets, commonly used in RNS keyswitching
        \item BV family
        \begin{itemize}
            \item HPS, closely related to BEHZ
            \item BV per limb 
        \end{itemize}
        \item GHS
        \item Hybrid
    \end{itemize}
}\fi

\subsection{Brakerski-Vaikuntanathan}
\begin{itemize}
    \item Digit decomposition is common in keyswitching and other applications, including \gls{gsw} ciphertexts and TFHE \cite{cryptoeprint:2012/099, cryptoeprint:2018/421, cryptoeprint:2011/277}
    \item Used extensively in \cite{BV}; decomposition-only key switching called BV-style
    \item Not directly applicable in RNS without multi-precision integers, but some 'translations' exist
\end{itemize}

\subsubsection{HPS Gadget}
\begin{itemize}
    \item First instantiated in \cite{cryptoeprint:2016/510}
    \item Uses CRT instead of number decomposition
\end{itemize}

\begin{align}
    \label{eq:bv:d}
    D_{Q}({a}) &= \left(\imod{a\left(\frac{Q}{q_0}\right)^{-1}}{0},\dots,\imod{a\left(\frac{Q}{q_i}\right)^{-1}}{i}\right) \\
    \label{eq:bv:p}
    P_{Q}({a}) &= \left(\qmod{a\frac{Q}{q_0}},\dots,\qmod{a\frac{Q}{q_i}}\right) 
    % = \left([a\hat{q}_0^{-1}]_{q_0},\dots,[a\hat{q}_{k-1}^{-1}]_{q_{k-1}}\right)
\end{align}

\begin{lemma}
    \label{lemma:bv_gadget}
    The vectors \eqref{eq:bv:d} and \eqref{eq:bv:p} satisfy the gadget property.
    \begin{equation}
        \langle D_Q(a), P_Q(b)\rangle = a\cdot b \pmod{Q}
    \end{equation}
\end{lemma}

\begin{itemize}
    \item Let $\hat{q}_j = Q/q_j$
    \item Halevi, Polyakov and Shoup \cite{cryptoeprint:2018/117} showed that the $\hat{q_j}$ factors can be moved into $P_Q(a)$ without violating Lemma \ref{lemma:bv_gadget}
    \item This makes the decomposition free for an RNS polynomial, as it simply isolates each tower
\end{itemize}

\begin{align}
    \hps{D}(a) &= \left( \imod{a}{0}, \imod{a}{1}, \cdots, \imod{a}{i}\right)\\
    \hps{P}(a) &= \left( a\cdot\imod{\hat{q}_0^{-1}}{0} \cdot \qmod{\hat{q}_0}, \dots, a\cdot\imod{\hat{q}_\ell^{-1}}{i} \cdot \qmod{\hat{q}_\ell} \right)
\end{align}

\begin{lemma}
    The HPS gadget applied to an RNS (CRT) polynomial is free, and equivalent to isolating each limb.
    \begin{equation*}
        \text{Let } e_j = [\hat{q}_j^{-1}]_{q_j}\cdot\hat{q}_j \bmod{Q}, \text{ then }
        e_j \bmod{q_i} = \begin{cases}
            1 & i = j \\
            0 & i \not= j
        \end{cases}
    \end{equation*}
    \begin{equation*}
        \text{and } \hps{P}(a) = (a\cdot e_0, a\cdot e_1, \dots, a\cdot e_{k - 1})
    \end{equation*}
    \begin{proof}
        \begin{align*}
            i = j &\;\implies\; e_j \equiv \hat{q}_j^{-1}\cdot 0 \equiv 0 &\pmod{q_i} \\
            i \not= j &\;\implies\; e_j \equiv \hat{q}_j^{-1}\cdot \hat{q}_j \equiv 1 &\pmod{q_i}
        \end{align*}
    \end{proof}
\end{lemma}

\begin{itemize}
    \item Note: The $e_j$ factors are static and can be pre-calculated
\end{itemize}
 
% \begin{equation}
%     \text{When the tower index $i$ is } 
%     \begin{cases}
%         i \not= j,      & \hat{Q}_j \equiv Q/q_j \equiv  ({q_0\cdots q_{i}\cdots q_{k-1}})/q_j \equiv n\cdot q_i \equiv 0 \pmod{q_i} \\
%         i = j,  & \hat{Q}_j^{-1}\hat{Q}_j \equiv 1 \pmod{q_i} 
%     \end{cases}
% \end{equation}

\iffable{
    \color{red}
    \subsection{Brakerski-Vaikuntanathan}
    % --- Method in the abstract, modulus-agnostic ---
    \begin{itemize}
        \item BV \cite{BV11} = decomposition-only key switching: no modulus extension,
              everything stays at modulus $Q$
        \item Split ciphertext element into small digits $D(a)$; key encrypts $s$ scaled by
              matching constants $P(s)$; correctness via gadget property
              $\langle D(a), P(b)\rangle \equiv a\cdot b \pmod{Q}$ (ref.\ abstract gadget section)
        \item Noise governed by digit size: $v_{\mathrm{ks}} = \langle D(c_1), \vec{e}\,\rangle$,
              so $\|v_{\mathrm{ks}}\|_\infty$ scales with $\max_j \|D(c_1)_j\|_\infty$
        \item Classical instantiation: radix-$\omega$ digits of each coefficient,
              $D_\omega(a) = (a_0, \dots, a_{\ell_\omega-1})$ with $a = \sum_m a_m \omega^m$,
              $P_\omega(b) = (b, b\omega, \dots, b\omega^{\ell_\omega - 1})$, digits bounded by $\omega/2$
        \item Transition / the problem: radix decomposition reads off the \emph{positional}
              representation of a coefficient --- in RNS this is never materialized;
              coefficients exist only as residues mod $q_j$, and reconstructing the
              multiprecision value would defeat the purpose of RNS
    \end{itemize}
    
    \subsubsection{RNS Instantiation}
    \begin{itemize}
        \item Resolution: no reconstruction needed --- \emph{the CRT is itself a decomposition}
        \item Each residue $\imod{a}{j}$ is a word-sized digit; the CRT reconstruction formula
              \begin{equation*}
                  a \equiv \sum_j \imod{a\,\hat{q}_j^{-1}}{j} \cdot \hat{q}_j \pmod{Q},
                  \qquad \hat{q}_j = Q/q_j
              \end{equation*}
              expresses $a$ as a linear combination of digits with fixed constants ---
              exactly the shape of a gadget
        \item Formalizing yields the RNS gadget, first given in \cite{cryptoeprint:2016/510}:
    \end{itemize}
    \begin{align}
        \label{eq:bv:d}
        D_{Q}(a) &= \left(\imod{a\,\hat{q}_0^{-1}}{0},\; \dots,\; \imod{a\,\hat{q}_\ell^{-1}}{\ell}\right) \\
        \label{eq:bv:p}
        P_{Q}(a) &= \left(\qmod{a\,\hat{q}_0},\; \dots,\; \qmod{a\,\hat{q}_\ell}\right)
    \end{align}
    
    \begin{lemma}[Gadget property of the RNS gadget]
        \label{lemma:bv_gadget}
        $\langle D_Q(a), P_Q(b)\rangle \equiv a\cdot b \pmod{Q}$.
        % Proof device: CRT idempotents e_j = [q̂_j^{-1}]_{q_j} · q̂_j mod Q,
        % e_j ≡ 1 (mod q_j), e_j ≡ 0 (mod q_i) for i ≠ j, hence Σ_j e_j ≡ 1 (mod Q);
        % then ⟨D_Q(a), P_Q(b)⟩ ≡ ab · Σ_j e_j ≡ ab (mod Q).
    \end{lemma}
    
    \begin{itemize}
        \item Remark (invariance under folding): the lemma is invariant under moving
              per-coordinate constants between $D$ and $P$ --- the $j$-th summands differ
              by a multiple of $q_j \hat{q}_j = Q$
        \item HPS \cite{cryptoeprint:2018/117} exploit this: move the
              $\imod{\hat{q}_j^{-1}}{j}$ factors from $D_Q$ into $P_Q$:
    \end{itemize}
    \begin{align}
        \hps{D}(a) &= \left( \imod{a}{0},\; \imod{a}{1},\; \dots,\; \imod{a}{\ell}\right)\\
        \hps{P}(a) &= \left( a\cdot\imod{\hat{q}_0^{-1}}{0}\cdot\qmod{\hat{q}_0},\;
                       \dots,\; a\cdot\imod{\hat{q}_\ell^{-1}}{\ell}\cdot\qmod{\hat{q}_\ell} \right)
                    = \left(a\cdot e_0,\; \dots,\; a\cdot e_\ell\right)
    \end{align}
    \begin{itemize}
        \item $\hps{D}(a)$ is exactly the RNS representation of $a$: decomposition is
              \emph{free} for an RNS polynomial --- limb lookup, zero arithmetic
        \item The $e_j = \imod{\hat{q}_j^{-1}}{j}\cdot\hat{q}_j \bmod Q$ are the CRT
              idempotents (textbook CRT, same objects as in the proof of
              Lemma~\ref{lemma:bv_gadget}); static, precomputed into the key at keygen
        \item Noise unchanged by folding: digit bound $q_j/2$ in both forms
        \item Closing / forward pointer: digits bounded by $\tilde{q}/2$ with
              $\tilde{q} = \max_j q_j$, so noise scales with limb size;
              when smaller digits are needed, each limb --- now word-sized --- can itself
              be radix-decomposed (next section). Note: the CRT gadget is what makes
              radix decomposition applicable again
    \end{itemize}
    
    \subsubsection{Per-Limb Digit Decomposition}
    \begin{itemize}
        \item Goal: reduce digit size from $\tilde{q}/2$ to $\omega/2$ for radix $\omega < \tilde{q}$
        \item Key fact: radix-$\omega$ decomposition of a digit is an \emph{exact integer
              identity}, $d_j = \sum_m \omega^m d_{j,m}$ over $\mathbb{Z}$ --- so it
              composes with \emph{any} gadget whose digits are small integer representatives
    \end{itemize}
    
    \begin{lemma}[Gadget composition]
        \label{lemma:double_decomp}
        Let $(D, P)$ satisfy the gadget property with integer digit representatives
        $|D(a)_j| \le q_j/2$. Then
        \begin{align*}
            D'(a) &= \left(d_{j,m}\right)_{j,m},
            \qquad \text{where } D(a)_j = \sum_m \omega^m d_{j,m},\; |d_{j,m}| \le \omega/2 \\
            P'(b) &= \left(\omega^m \cdot P(b)_j\right)_{j,m}
        \end{align*}
        satisfies $\langle D'(a), P'(b)\rangle \equiv a\cdot b \pmod{Q}$.
        % Proof: substitute the radix expansion into ⟨D(a), P(b)⟩, apply linearity. Two lines.
    \end{lemma}
    
    \begin{itemize}
        \item Instantiation with the RNS gadget (HPS form): $D'(a)$ = radix decomposition
              of each limb of the RNS representation
              \begin{align*}
                  D'(a) &= \left( \left(\imod{a}{j}\right)_m \right)_{j,m}, \qquad
                  P'(b) = \left( b \cdot \omega^m e_j \right)_{j,m}
              \end{align*}
        \item Still no scalar multiplications on the ciphertext side: all constants
              ($e_j$, $\omega^m$) live in $P'$, precomputable; ciphertext side is limb
              lookup + word-sized radix split
        \item Cost/noise tradeoff: $\ell_{\omega,\tilde{q}} = \lfloor\log_\omega \tilde{q}\rfloor + 1$
              times more gadget components (key size, NTTs), digits shrink
              $\tilde{q}/2 \to \omega/2$; noise bound as in \cite{cryptoeprint:2021/204},
              App.~B.2.1 (second level of decomposition)
        \item Naming remark: no modulus extension is performed, so this remains a BV-type
              key switch; we refer to it simply as \emph{BV} in the remainder of this
              thesis (and in the implementation)
        \item Remark (forward pointer): Lemma~\ref{lemma:double_decomp} applies to any
              gadget with small integer digits --- in particular to the digits of hybrid
              key switching (Section~\ref{sec:hybrid})
    \end{itemize}
}\fi


% \newpage
% \iffable{
%     \color{purple}

%     \textbf{The "broadcast" is not an operation — it's just what a small integer looks like in RNS.}
%     A digit $d$ is a small signed integer, $|d| \le B/2$. Your gadget theorem needs $D(c)$ to be a vector of small ring elements in $R_Q$, because each digit polynomial gets multiplied against a full RGSW row (whose mask and noise occupy all $k$ towers — only the $m\cdot G$ part is sparse). And the RNS representation of the integer $d$ is, by definition,
    
%     $$d \;\longmapsto\; \big([d]{q_0}, [d]{q_1}, \dots, [d]{q{k-1}}\big),$$
    
%     which for $d \ge 0$ is literally the number $d$ written $k$ times, and for $d < 0$ is $q_t - |d|$ in each tower. That's the entire "broadcast": storing a small integer in RNS form. In a non-RNS implementation this step wouldn't even be a line of code — a digit would just be a ring element. It only looks like a named, load-bearing operation in context-bv.cpp:245-250 because OpenFHE's data layout forces you to materialize every tower explicitly, and because signed digits need the per-tower negative encoding. In the thesis it deserves one sentence, not a name: "each digit $d_{j,i}$ is a small integer; we lift it to $R_Q$ by reducing it modulo each $q_t$." I'd drop the word "broadcast" entirely from both the write-up and, if you feel like it, the code comments.
    
%     So Decompose is conceptually just two steps, one interesting and one trivial:
    
%     Per-tower digit extraction (the interesting one): for each tower $j$, coefficient-wise signed base-$B$ decomposition of the single-word representative of $[c]_{q_j}$. No multiprecision, no CRT reconstruction — this is the payoff of the whole construction.
%     Lift to $R_Q$ (the trivial one): each digit, being a small integer, is written into all towers as above, then NTT'd. The only real cost here is the $k$ NTTs per digit polynomial — that's the honest bottleneck of the external product, worth one remark.
    
%     \textbf{The one sentence that makes the whole optimization constructive.}
%     Here is the framing I'd hang everything on, because it turns the "0-towers are skipped" hand-wave into a precise statement:
    
%     In RNS layout, multiplying by $e_j = [\hat q_j^{-1}]_{q_j}\hat q_j$ is not a computation — it is array indexing. Since $a\cdot e_j$ is zero in every tower except tower $j$ (where it equals $[a]_{q_j}$), "multiply by $e_j$" means "select row $j$ of the limb matrix."
    
%     Look at your code through that lens and notice: the idempotents $e_j$ never appear anywhere in it. PowersOfBase doesn't compute $m B^i e_j$ and then discard zero towers — it computes only the $k\ell$ limbs $[mB^i]{q_j}$ that are nonzero (context-bv.cpp:184-188), and EncryptRGSW performs the "multiplication by $e_j$" by indexing tower $j$ (context-bv.cpp:115-117). Likewise the layer-1 decomposition $D^{(1)}(c) = ([c]{q_0},\dots,[c]{q{k-1}})$ costs nothing — it's reading the limbs you already have — which is why Decompose jumps straight to digit extraction per tower. The optimized algorithms are what you get when every multiplication by $e_j$ and every evaluation of $D^{(1)}$ is replaced by indexing. That is the constructive story, and it's exactly what makes the pseudocode fall out.
    
%     (One correction to your framing: the zero-tower skipping lives on the $P$/gadget side — PowersOfBase and EncryptRGSW. In Decompose nothing is zero; what's exploited there is input locality: the $\ell$ digit polynomials of group $j$ read only limb $j$. Both are consequences of the same sparsity fact, but they're different claims, and stating them separately will make the algorithms easier to justify.)
% }\fi


\newpage
\chapter{Implementation}
\section{Brakerski-Vaikuntanathan}
As established, digit-decomposing each RNS prime with $\ell$ digits is equivalent to digit-decomposing the full polynomial with $k\cdot\ell$ digits, where $Q = \prod_{i=0}^{k-1}q_i$ or the basis $\mathcal{B}_{Q}$.

\subsection{Powers of Base}
\begin{enumerate}
    \item $a \to (a, aB, \dots, B^{\ell-1})$
    \item $g = (1, B, \cdots, B^{\ell}, B, \dots, B, B^2, \dots, B^{\ell-1}, \dots B^{\ell-1})$ ??
\end{enumerate}

\begin{algorithm}
\caption{Powers of Base (Optimized)}
\label{alg:bv:p}
\begin{algorithmic}[1]
\Require RNS coefficients $x = ([x]_{q_0}, ..., [x]_{q_{k-1}})$, base $B = 2^{b}$, number of digits $\ell$
\Ensure $\ell$ RNS/CRT polynomials $(a, aB, \dots, B^{\ell-1})$
\State Initialize a vector $x'$ of size $\ell$
\For{$i \in [0, \ell)$ \textbf{in parallel}}
    \For{$j \in [0, k)$}
        \State $x'_i[j] := x[j] \cdot B^{i} \bmod{q_j}$
    \EndFor
\EndFor
\State \Return $x'$
\end{algorithmic}
\end{algorithm}

\begin{itemize}
    \item $B$ (or $\ell$) is implicit from $\ell = \lfloor\log_{B}{(|\max_i{q_i}|)}\rfloor + 1$, and since $B = 2^b$ is a power of 2 we get $\ell = \lfloor \log(\max_i{q_i})/b)\rfloor + 1$.
    \item $B^i$ or $B^i \bmod{q_j}$ can be calculated ahead of time. 
    \item The for loop can be flattened to $0 \leq idx < k\cdot \ell$ to allow more parallelization.
    \item For maximum efficiency, the function can be encoded into the encryption function.
    \item The fully optimized encryption functions becomes:
\end{itemize}

\begin{algorithm}
\caption{Encrypt BV-RGSW (Optimized)}
\label{alg:bv:encrypt}
\begin{algorithmic}[1]
\Require Plaintext $m = ([m]_{q_0}, ..., [m]_{q_{k-1}})$, base $B = 2^{b}$, number of digits $\ell$
\Ensure $\rgsw(m)$ with $2k\ell$ rows
\State Initialize a vector $C$ of $2k\ell$ \texttt{DCRTPoly} objects $(c_0, c_1)$
\State Set $m$ to RNS-coefficients mode
\For{$idx \in [0, \ell\cdot k)$ \textbf{in parallel}}
    \State $(i, j) \gets [\lfloor idx / \ell \rfloor, \; idx \bmod{\ell}]$
    \State $mg \gets m[j] \cdot B^{i} \bmod{q_j}$
    \Comment{Retrieve $B^i\bmod{q_j}$ from lookup table}
    \State $C_{idx} \gets \mathsf{HE.Encrypt}(0) + (\mathsf{NTT}(mg), 0)$
    \State $C_{idx + \ell} \gets \mathsf{HE.Encrypt}(0) + (0, \mathsf{NTT}(mg))$
\EndFor
\State \Return $C$
\end{algorithmic}
\end{algorithm}

% \newpage

% \begin{minted}[bgcolor=gray]{c++}
% std::vector<Poly> ExtendedContextBVImpl::PowersOfBase(const Poly& input) const 
% {
%     const auto n_towers = this->GetElementParams()->GetParams().size();

%     std::vector<Poly> result(m_ell);
%     result[0] = input;

% #pragma omp parallel for 
%     for (size_t i = 1; i < m_ell; i++) {
%         Poly scaled(input.GetParams(), input.GetFormat(), true);

%         for (size_t j = 0; j < n_towers; j++) {
%             auto factor = GetPower(i, j);
%             auto limb = input.GetElementAtIndex(j).Times(factor);
%             scaled.SetElementAtIndex(j, std::move(limb));
%         }

%         result[i] = std::move(scaled);
%     }

%     return result;
% }
% \end{minted}

% \begin{minted}[bgcolor=gray]{c++}
% RGSW ExtendedContextBVImpl::EncryptRGSW(const PublicKey& pk, const Plaintext& pt) const 
% {
%     // clang-format off
%     const auto msg = pt->GetElement<Poly>();
%     const auto zero = IsCoefPackedPlaintext(pt)
%         ? this->MakeCoefPackedPlaintext({0})
%         : this->MakePackedPlaintext({0});
%     // clang-format on

%     std::vector<Poly> digits = PowersOfBase(msg);

%     const size_t k = msg.GetNumOfElements();
%     const size_t half = k * m_ell;
%     const size_t full = 2 * half;

%     RGSW rows(full);

%     // For noiseless mode, encrypt zero once so we can clone its metadata
%     RLWE noiseless_template;
%     if (!noisy) {
%         noiseless_template = this->Encrypt(pk, zero);
%     }

% #pragma omp parallel for num_threads(lbcrypto::OpenFHEParallelControls.GetThreadLimit(full))
%     for (size_t row = 0; row < full; row++) {
%         size_t l = row / half;    // 0 for Upper (c_0), 1 for Lower (c_1)
%         size_t rem = row % half;  // The flat index within the current half
%         size_t j = rem / m_ell;   // The target tower [0 to k-1]
%         size_t i = rem % m_ell;   // The base power   [0 to ell-1]

%         RLWE ct;
%         if (!noisy) {
%             ct = noiseless_template->CloneEmpty();
%             Poly c0(this->GetElementParams(), Format::EVALUATION, true);
%             Poly c1(this->GetElementParams(), Format::EVALUATION, true);
%             ct->SetElements({std::move(c0), std::move(c1)});
%         } else {
%             ct = this->Encrypt(pk, zero);
%         }
%         auto elements = ct->GetElements();

%         auto target_limb = elements[l].GetElementAtIndex(j);
%         target_limb += digits[i].GetElementAtIndex(j);
%         elements[l].SetElementAtIndex(j, std::move(target_limb));
%         ct->SetElements(std::move(elements));
%         rows[row] = std::move(ct);
%     }

%     return rows;
% }
% \end{minted}

\newpage

\subsection{Decompose}
\begin{enumerate}
    \item 
\end{enumerate}

\begin{algorithm}
\caption{BV Decompose (Optimized)}
\label{alg:bv:d}
\begin{algorithmic}[1]
\Require Coefficients $x = ([x]_{q_0}, ..., [x]_{q_{k-1}})$, base $B = 2^b$, number of digits $\ell$
\Ensure $\ell$ polynomials of signed digits $(d_0, d_1, \dots, d_{\ell-1})$ where $d_j \in [-B/2, B/2)$
\State $a' \gets a$
\For{xxxx}
    \State xxxx
    
\EndFor
\State \Return $\{d\}$
\end{algorithmic}
\end{algorithm}

\begin{itemize}
    \item Optimized/folded form:
\end{itemize}

\begin{algorithm}
\caption{External Product $x \boxdot Y$ (Optimized BV-RGSW)}
\label{alg:bv:extprod}
\begin{algorithmic}[1]
    \Require RLWE ciphertext $x = (x_0, x_1) \in R_q^2$, RGSW ciphertext $Y=\rgsw(m)$ with $2k\ell$ rows, base $B = 2^b$, decomposition parameter/digits $\ell$
    \Ensure RLWE ciphertext encrypting $m \cdot m_x$
    
    \State $\mathsf{acc} \gets (0, 0)$
    \Comment{accumulator in evaluation (NTT) form}

    \For{$c\in \{0, 1\}$}
        \Comment{decompose each ciphertext component}
        \For{$idx \in [0, \ell \cdot k)$ \textbf{in parallel}}
            \State $(i, j) \gets [\lfloor idx/\ell \rfloor,\; idx \bmod \ell]$
        
            \State $r \gets [x_c]_{q_j}$
            \Comment{$j$-th CRT residue, coefficient form}
            
            \State $d_{idx} \gets \mathsf{Digit}_B(r, i)$
            \Comment{$i$-th base-$B$ digit; coeffs in $[0,B)$ or $(-B/2,B/2]$}

            \State $\tilde{d}_{idx} \gets \mathsf{CRTExtend}(d_{idx})$
            \Comment{lift small digit poly to all $k$ moduli, then NTT}
            
        \EndFor
        \State $\mathsf{acc} \gets \mathsf{acc} + \sum_{idx} \tilde{d}_{idx} \cdot Y_{idx + c \cdot k\ell}$
        \Comment{2 poly mults per row}
    \EndFor
    \State \Return $\mathsf{acc}$
\end{algorithmic}
\end{algorithm}